\documentclass[12pt,a4paper]{ctexart}

\usepackage{graphicx}
\usepackage{float}
\usepackage{geometry}
\usepackage{fancyhdr}
\usepackage{booktabs}
\usepackage{hyperref}

\geometry{left=2.5cm,right=2.5cm,top=2.5cm,bottom=2.5cm}

\pagestyle{fancy}
\fancyhf{}
\fancyhead[L]{\small 转转——校园二手物品交易平台}
\fancyhead[R]{\small 张可凡·全栈开发工作汇报}
\fancyfoot[C]{\thepage}

\title{{\Huge 全栈开发工作汇报}\\[0.3cm]{\Large 校园二手物品交易平台——转转}}
\author{{\large 张可凡（队长）}\\{\small 全栈开发 / 项目管理 / 测试}}
\date{2026年7月4日 — 2026年7月10日}

\begin{document}
\maketitle
\thispagestyle{fancy}

\section{个人角色与职责}

\begin{table}[H]
\centering
\begin{tabular}{lp{10cm}}
\toprule
\textbf{项目} & \textbf{内容} \\
\midrule
姓名 & 张可凡 \\
角色 & 队长 / 全栈开发 / 项目管理 / 测试 \\
主要职责 & Docker 容器化部署、管理后台前端开发、Postman 接口测试、文档整合 \\
Git 提交 & 18 次 \\
新增代码 & $\sim$2500 行 \\
\bottomrule
\end{tabular}
\end{table}

\section{每日工作内容}

\subsection{7月4日（周六）—— 项目启动与环境搭建}
\begin{itemize}
    \item 创建 GitHub 仓库并初始化目录结构
    \item 编写 \texttt{docker-compose.yml} 编排 MySQL + Redis + 后端 + 前端容器
    \item 配置 Nginx
    \item 编写 \texttt{README.md}
\end{itemize}

\subsection{7月5日（周日）—— 基础框架与核心后端}
\begin{itemize}
    \item 完成 Docker 全栈环境联调验证（phpMyAdmin + 后端 + 前端）
    \item 编写 Redis 配置与缓存工具类
    \item 编写 \texttt{deploy/mysql/my.cnf}、\texttt{deploy/redis/redis.conf}
\end{itemize}

\subsection{7月6日（周一）—— 用户模块 + 商品模块}
\begin{itemize}
    \item 编写需求分析文档（软件需求规格说明书）
    \item 编写系统架构设计文档
    \item 编写 \texttt{deploy/mysql/migration-indexes.sql} 索引优化脚本
\end{itemize}

\subsection{7月7日（周二）—— 订单模块 + 商品管理}
\begin{itemize}
    \item 完成数据库设计说明书
    \item 开始编写接口设计说明书
    \item 开发管理后台用户管理页（列表 + 搜索 + 禁用/启用）
\end{itemize}

\subsection{7月8日（周三）—— 消息模块 + 后台管理}
\begin{itemize}
    \item 完成接口设计说明书
    \item 开发管理后台商品审核页（待审核列表 + 通过/驳回/违规下架）
    \item 开发管理后台订单管理页
\end{itemize}

\subsection{7月9日（周四）—— 管理后台 + 联调}
\begin{itemize}
    \item 开发管理后台公告管理页（CRUD + 列表）
    \item 开发管理后台数据统计页（ECharts 图表）
    \item 使用 Postman 进行全量接口测试
\end{itemize}

\subsection{7月10日（周五）—— 收尾与文档}
\begin{itemize}
    \item 编写测试报告文档
    \item 编写用户使用说明书
    \item Postman 回归测试确认全部接口通过
    \item 汇总全组文档并校验格式
    \item 编写项目 README 最终版
\end{itemize}

\end{document}
